% Source problem: https://www.ap.t.u-tokyo.ac.jp/themes/ap_t_u-tokyo/graduate/pdf/2024ap_exam.pdf
\documentclass[11pt,a4paper]{article}
\usepackage[margin=24mm]{geometry}
\usepackage{fontspec}
\usepackage{xeCJK}
\setCJKmainfont{HaranoAjiMincho-Regular.otf}
\setCJKsansfont{HaranoAjiGothic-Medium.otf}
\usepackage{amsmath,amssymb}
\usepackage{tikz}
\setlength{\parindent}{1em}
\setlength{\parskip}{0.35em}
\title{物理工学専攻 院試解答集\\2024年度 物理学 第3問}
\author{政岡凜太郎}
\date{}
\begin{document}
\maketitle

以下で $C=(\partial U/\partial T)$ は指定された振動子系の全熱容量とする。固体の単位質量あたりの格子比熱は、得られた $C$ を固体の質量で割ったものである。体積や振動数は温度に依存しないものとして微分する。

\section*{[1] 古典的な1次元調和振動子}
\subsection*{[1.1] 分配関数}
$M$ 個の振動子は独立なので、まず1個の位相空間積分 $z_{\mathrm{cl}}$ を計算する。運動量と位置の積分を分け、問題文のガウス積分を使うと
\begin{align}
z_{\mathrm{cl}}
&=\frac{1}{2\pi\hbar}
 \int_{-\infty}^{\infty}\!\mathrm dp\,
 e^{-\beta p^2/(2m)}
 \int_{-\infty}^{\infty}\!\mathrm dx\,
 e^{-\beta m\omega^2x^2/2}\\
&=\frac{1}{2\pi\hbar}
 \sqrt{\frac{2\pi m}{\beta}}
 \sqrt{\frac{2\pi}{\beta m\omega^2}}
 =\frac{1}{\beta\hbar\omega}.
\end{align}
従って
\begin{equation}
\boxed{Z(T)=z_{\mathrm{cl}}^M
 =\left(\frac{k_{\mathrm B}T}{\hbar\omega}\right)^M}.
\end{equation}

\subsection*{[1.2] 内部エネルギー}
正準分布の関係 $U=-\partial\ln Z/\partial\beta$ に上式を代入する。$\ln Z=-M\ln(\beta\hbar\omega)$ より
\begin{equation}
\boxed{U=\frac{M}{\beta}=Mk_{\mathrm B}T}.
\end{equation}
1振動子に運動・位置の二つの二次項があり、各項が $k_{\mathrm B}T/2$ を担うという等分配則とも一致する。

\subsection*{[1.3] 熱容量}
振動数を固定して $U$ を温度微分すると
\begin{equation}
\boxed{C=\left(\frac{\partial U}{\partial T}\right)_{\omega}=Mk_{\mathrm B}}.
\end{equation}

\section*{[2] 単一振動数の量子的調和振動子}
\subsection*{[2.1] 分配関数}
1振動子について $n=0,1,\ldots$ のボルツマン因子を和にする。$q=e^{-\beta\hbar\omega}$ は $0<q<1$ なので等比級数を使え、
\begin{equation}
z_{\mathrm q}=\sum_{n=0}^{\infty}
 e^{-\beta\hbar\omega(n+1/2)}
 =\frac{e^{-\beta\hbar\omega/2}}{1-e^{-\beta\hbar\omega}}
 =\frac{1}{2\sinh(\beta\hbar\omega/2)}.
\end{equation}
$M$ 個の独立なモードでは
\begin{equation}
\boxed{Z(T)=\left[
\frac{e^{-\beta\hbar\omega/2}}{1-e^{-\beta\hbar\omega}}
\right]^M}.
\end{equation}

\subsection*{[2.2] 内部エネルギー}
$\ln Z=M[-\beta\hbar\omega/2-\ln(1-e^{-\beta\hbar\omega})]$ を $\beta$ で微分すると
\begin{equation}
\boxed{U=M\hbar\omega\left(
\frac12+\frac{1}{e^{\beta\hbar\omega}-1}
\right)}.
\end{equation}
第1項は $T=0$ でも残る零点エネルギー、第2項は熱励起のエネルギーである。

\subsection*{[2.3] 熱容量}
$\Theta=\hbar\omega/k_{\mathrm B}$ を振動子の温度尺度とおくと、$\beta\hbar\omega=\Theta/T$ である。零点項は温度微分で消え、$y=\Theta/T$ について $\mathrm dy/\mathrm dT=-y/T$ を使えば
\begin{equation}
\boxed{C=Mk_{\mathrm B}
\left(\frac{\Theta}{T}\right)^2
\frac{e^{\Theta/T}}{(e^{\Theta/T}-1)^2}}
=Mk_{\mathrm B}\left[
\frac{\Theta/(2T)}{\sinh(\Theta/(2T))}
\right]^2.
\end{equation}

\subsection*{[2.4] 低温・高温と概形}
$T\ll\Theta$ では $e^{\Theta/T}\gg1$、$T\gg\Theta$ では $\sinh(y/2)\sim y/2$ である。従って
\begin{equation}
\boxed{C\sim Mk_{\mathrm B}\left(\frac{\Theta}{T}\right)^2e^{-\Theta/T}
\quad(T\ll\Theta)},\qquad
\boxed{C\longrightarrow Mk_{\mathrm B}\quad(T\gg\Theta)}.
\end{equation}
低温では最初の励起に必要な有限のエネルギー $\hbar\omega$ が熱励起を抑える。図の横軸 $T/\Theta$、縦軸 $C/(Mk_{\mathrm B})$ は無次元である。
\begin{center}
\begin{tikzpicture}[x=3.0cm,y=3.0cm]
 \draw[->] (0,0)--(3.2,0) node[right] {$T/\Theta$};
 \draw[->] (0,0)--(0,1.2) node[above] {$C/(Mk_{\mathrm B})$};
 \draw[gray,densely dotted] (0,1)--(3,1);
 \node[left] at (0,1) {$1$};
 \foreach \x in {1,2,3} {\draw (\x,0)--(\x,-0.025) node[below] {$\x$};}
 \draw[blue,very thick] plot[smooth] coordinates {
 (0,0) (0.05,0.000001) (0.1,0.004540) (0.15,0.056706)
 (0.2,0.170742) (0.3,0.426247) (0.4,0.608890)
 (0.5,0.724062) (0.7,0.845975) (1,0.920674)
 (1.5,0.963772) (2,0.979425) (3,0.990792)};
 \node[blue] at (1.2,0.70) {量子振動子};
\end{tikzpicture}
\end{center}

\section*{[3] デバイ型の振動数分布}
\subsection*{[3.1] 内部エネルギーと熱容量の積分}
まず与えられた状態密度の規格化は
\begin{equation}
\int_0^{\omega_{\mathrm D}}g(\omega)\,\mathrm d\omega
=\frac{9N}{\omega_{\mathrm D}^3}
\int_0^{\omega_{\mathrm D}}\omega^2\,\mathrm d\omega=3N
\end{equation}
であり、3次元固体の振動モード数に一致する。各周波数のモードに[2.2]の1振動子のエネルギーを適用して、
\begin{equation}
\boxed{U=\int_0^{\omega_{\mathrm D}}\!g(\omega)\hbar\omega
\left[\frac12+\frac{1}{e^{\beta\hbar\omega}-1}\right]
\mathrm d\omega}.
\end{equation}
このうち零点エネルギーは $\frac12\int g(\omega)\hbar\omega\,\mathrm d\omega=9N\hbar\omega_{\mathrm D}/8$ であり、熱容量に寄与しない。従って、微分を積分の中に入れると
\begin{equation}
\boxed{C=k_{\mathrm B}\int_0^{\omega_{\mathrm D}}\!g(\omega)
(\beta\hbar\omega)^2
\frac{e^{\beta\hbar\omega}}{(e^{\beta\hbar\omega}-1)^2}
\mathrm d\omega}.
\end{equation}

\subsection*{[3.2] 低温・高温と概形}
デバイ温度 $\Theta_{\mathrm D}=\hbar\omega_{\mathrm D}/k_{\mathrm B}$ を導入し、積分変数を $x=\beta\hbar\omega$ に変える。上限は $\Theta_{\mathrm D}/T$、また $g(\omega)\mathrm d\omega=9N(T/\Theta_{\mathrm D})^3x^2\mathrm dx$ なので
\begin{equation}
C=9Nk_{\mathrm B}\left(\frac{T}{\Theta_{\mathrm D}}\right)^3
\int_0^{\Theta_{\mathrm D}/T}
\frac{x^4e^x}{(e^x-1)^2}\,\mathrm dx.
\end{equation}
$T\ll\Theta_{\mathrm D}$ では上限が十分大きい。被積分関数は大きい$x$で$x^4e^{-x}$に近づくため、上限から先の尾部は指数的に小さく、積分の上限を無限大に延ばせる。問題文の積分値を使い、
\begin{equation}
\boxed{C\sim\frac{12\pi^4}{5}Nk_{\mathrm B}
\left(\frac{T}{\Theta_{\mathrm D}}\right)^3
\quad(T\ll\Theta_{\mathrm D})}.
\end{equation}
高温では全モードに $k_{\mathrm B}$ の熱容量が戻る。実際、$x\ll1$ で $x^2e^x/(e^x-1)^2\to1$ だから、規格化したモード数を用いて
\begin{equation}
\boxed{C\longrightarrow 3Nk_{\mathrm B}\quad(T\gg\Theta_{\mathrm D})}.
\end{equation}
概形は低温で原点から $T^3$ に沿って立ち上がり、上側の漸近値 $3Nk_{\mathrm B}$ に近づく。
\begin{center}
\begin{tikzpicture}[x=3.0cm,y=3.0cm]
 \draw[->] (0,0)--(3.2,0) node[right] {$T/\Theta_{\mathrm D}$};
 \draw[->] (0,0)--(0,1.2) node[above] {$C/(3Nk_{\mathrm B})$};
 \draw[gray,densely dotted] (0,1)--(3,1);
 \node[left] at (0,1) {$1$};
 \foreach \x in {1,2,3} {\draw (\x,0)--(\x,-0.025) node[below] {$\x$};}
 \draw[blue,very thick] plot[smooth] coordinates {
 (0,0) (0.05,0.009741) (0.1,0.075821) (0.15,0.212992)
 (0.2,0.368635) (0.3,0.607703) (0.4,0.745853)
 (0.5,0.825408) (0.7,0.904954) (1,0.951732)
 (1.5,0.978126) (2,0.987611) (3,0.994466)};
 \node[blue] at (1.05,0.73) {デバイ模型};
\end{tikzpicture}
\end{center}

\section*{[4] 3種類の固体モデルの比較}
\subsection*{[4.1] 同じ座標でのプロット}
[1]と[2]の $M$ を $3N$ に替えると、$\Theta=\hbar\omega/k_{\mathrm B}$ を用いて
\begin{align}
C_{\mathrm I}&=3Nk_{\mathrm B},\\
C_{\mathrm{II}}&=3Nk_{\mathrm B}
\left(\frac{\Theta}{T}\right)^2
\frac{e^{\Theta/T}}{(e^{\Theta/T}-1)^2},\\
C_{\mathrm{III}}&=9Nk_{\mathrm B}
\left(\frac{T}{\Theta_{\mathrm D}}\right)^3
\int_0^{\Theta_{\mathrm D}/T}
\frac{x^4e^x}{(e^x-1)^2}\,\mathrm dx.
\end{align}
原題は $\omega$ と $\omega_{\mathrm D}$ の比を指定していない。以下の図は比較例として $\omega=\omega_{\mathrm D}$、すなわち $\Theta=\Theta_{\mathrm D}$ と置き、縦軸を $3Nk_{\mathrm B}$ で割った。
\begin{center}
\begin{tikzpicture}[x=3.0cm,y=3.2cm]
 \draw[->] (0,0)--(3.2,0) node[right] {$T/\Theta$};
 \draw[->] (0,0)--(0,1.19) node[above] {$C/(3Nk_{\mathrm B})$};
 \node[left] at (0,1) {$1$};
 \foreach \x in {1,2,3} {\draw (\x,0)--(\x,-0.023) node[below] {$\x$};}
 \draw[black,thick,densely dotted] (0,1)--(3,1);
 \draw[blue,very thick] plot[smooth] coordinates {
 (0,0) (0.05,0.009741) (0.1,0.075821) (0.15,0.212992)
 (0.2,0.368635) (0.3,0.607703) (0.4,0.745853)
 (0.5,0.825408) (0.7,0.904954) (1,0.951732)
 (1.5,0.978126) (2,0.987611) (3,0.994466)};
 \draw[red,very thick,dashed] plot[smooth] coordinates {
 (0,0) (0.05,0.000001) (0.1,0.004540) (0.15,0.056706)
 (0.2,0.170742) (0.3,0.426247) (0.4,0.608890)
 (0.5,0.724062) (0.7,0.845975) (1,0.920674)
 (1.5,0.963772) (2,0.979425) (3,0.990792)};
 \draw[black,thick,densely dotted] (0.20,-0.21)--(0.58,-0.21);
 \node[anchor=west] at (0.62,-0.21) {$C_{\mathrm I}$ 古典};
 \draw[red,very thick,dashed] (1.18,-0.21)--(1.56,-0.21);
 \node[anchor=west] at (1.60,-0.21) {$C_{\mathrm{II}}$ 単一振動数};
 \draw[blue,very thick] (0.20,-0.39)--(0.58,-0.39);
 \node[anchor=west] at (0.62,-0.39) {$C_{\mathrm{III}}$ デバイ模型};
\end{tikzpicture}
\end{center}

\subsection*{[4.2] 高温での比較}
$T\gg\Theta,\Theta_{\mathrm D}$ では各量子モードの平均エネルギーの主項が $k_{\mathrm B}T$ となる。従って $3N$ 個の振動モードはいずれも等分配則の値に戻り、
\begin{equation}
\boxed{C_{\mathrm I}=3Nk_{\mathrm B},\qquad
C_{\mathrm{II}}\longrightarrow3Nk_{\mathrm B},\qquad
C_{\mathrm{III}}\longrightarrow3Nk_{\mathrm B}}.
\end{equation}
有限の高温での近づき方も確認できる。$y^2e^y/(e^y-1)^2=1-y^2/12+O(y^4)$ を各モードに用いると
\begin{equation}
C_{\mathrm{II}}=3Nk_{\mathrm B}
\left[1-\frac1{12}\left(\frac{\Theta}{T}\right)^2+O\!\left((\Theta/T)^4\right)\right],\qquad
C_{\mathrm{III}}=3Nk_{\mathrm B}
\left[1-\frac1{20}\left(\frac{\Theta_{\mathrm D}}{T}\right)^2+O\!\left((\Theta_{\mathrm D}/T)^4\right)\right].
\end{equation}
後者では$g(\omega)$による平均が$\langle\omega^2\rangle=(3N)^{-1}\int_0^{\omega_{\mathrm D}}g(\omega)\omega^2\,d\omega=3\omega_{\mathrm D}^2/5$であり、$1/12$にこれを掛けて係数$1/20$を得る。どちらの量子曲線も古典値へ下から近づく。

\subsection*{[4.3] 低温での比較}
古典模型を形式的に低温へ外挿すると $C_{\mathrm I}=3Nk_{\mathrm B}$ のままであり、熱励起の量子化による凍結を表せない。一方、単一振動数の量子模型にはすべてのモードに同じ有限の励起ギャップ $\hbar\omega$ があるので
\begin{equation}
\boxed{C_{\mathrm{II}}\sim3Nk_{\mathrm B}
\left(\frac{\Theta}{T}\right)^2e^{-\Theta/T}}.
\end{equation}
デバイ模型には $\omega\to0$ の低周波モードが連続して存在する。温度 $T$ で励起しやすい $\hbar\omega\lesssim k_{\mathrm B}T$ のモード数は $g(\omega)\propto\omega^2$ から $T^3$ に比例するため、
\begin{equation}
\boxed{C_{\mathrm{III}}\sim\frac{12\pi^4}{5}Nk_{\mathrm B}
\left(\frac{T}{\Theta_{\mathrm D}}\right)^3}.
\end{equation}
したがって $T\to0$ では両量子模型の熱容量は零になるが、デバイ模型の冪減衰は単一振動数模型の指数減衰より遅い。ここでは格子振動だけを比較している。実在の金属では電子の熱容量が低温でおおむね $T$ に比例する場合があり、この3模型の式には含まれない。

\end{document}
